\documentclass{report}
\usepackage[utf8]{inputenc}
\usepackage[T1]{fontenc}
\usepackage{geometry}
\usepackage{graphicx}
\usepackage{float}
\usepackage{hyperref}
\usepackage[english,indonesian]{babel}

\title{Optimasi Pemrosesan Template LaTeX Menggunakan Pendekatan Berbasis AI}
\author{Academic Support Agent}
\date{\today}

\begin{document}
\maketitle

\begin{abstract}

\end{abstract}

\chapter{PENDAHULUAN}


\chapter{TINJAUAN PUSTAKA}


% Original content (commented out):
% \documentclass[12pt, a4paper,twoside]{report}

%% Every LaTeX document begins with a preamble, which loads packages and
%% defines various settings to make the document look right. Mostly,
%% you can ignore everything in this template before \begin{document} on line 74

\usepackage{mathtools,amsthm,amsfonts} % Enable useful mathematical symbols/environments
\usepackage{graphicx} % Enable graphics
\usepackage{fancyhdr,titlesec,microtype} % enable various formatting commands
\usepackage[margin=2.5cm]{geometry} % Set margin size
\usepackage{palatino} % Set the font
\usepackage[latin1]{inputenc} % Allow you to input accents, umlauts and other characters
\usepackage[T1]{fontenc} % Lets LaTeX print a wider array of characters
\usepackage{tikz,tikz-3dplot,tkz-euclide} % Enable tikz drawings

\usepackage{xcolor} % Enable coloured elements
\definecolor{mypurple}{HTML}{622567} %%% Purple
\definecolor{myred}{HTML}{D55C19} %%%EssexOrange
\definecolor{myblue}{HTML}{007A87} %%%Seagrass

% For technical reasons, hyperref should be loaded after all other packages
\usepackage[colorlinks,linkcolor=myblue,citecolor=mypurple]{hyperref}

\renewcommand{\baselinestretch}{1.5} % 1.5 line spacing

% Define \begin{theorem}, \end{theorem}, etc.
\theoremstyle{plain} % The following will be italicised
\newtheorem{theorem}{Theorem}[chapter]
\newtheorem{lemma}[theorem]{Lemma}
\newtheorem{proposition}[theorem]{Proposition}
\newtheorem{corollary}[theorem]{Corollary}

\theoremstyle{definition} % The following environments will not use italics
\newtheorem{definition}[theorem]{Definition}
\newtheorem{example}[theorem]{Example}

\theoremstyle{remark} % The following environments will not use italics or bold titles
\newtheorem{remark}[theorem]{Remark}

\numberwithin{equation}{chapter}

% Fancy headings
\pagestyle{fancy}
\setlength{\headheight}{15pt}
\fancyheadoffset[LE,RO]{0pt}
\renewcommand{\chaptermark}[1]{\markboth{#1}{}}
\renewcommand{\sectionmark}[1]{\markright{\thesection\ #1}}
\fancyhf{}
\fancyhead[LE]{\makebox[0pt][l]{\thepage}\hfill\leftmark}
\fancyhead[RO]{\rightmark\hfill\makebox[0pt][r]{\thepage}}
\fancypagestyle{plain}{%
    \fancyhead{} % get rid of headers
    \renewcommand{\headrulewidth}{0pt} % and the line
}

% Fancy chapter numbers
\titleformat{\chapter}[display]
    {\normalfont\bfseries\color{myred}}
    {\filleft\hspace*{-60pt}%
        \rotatebox[origin=c]{90}{%
            \normalfont\color{black}\Large%
            \textls[180]{\textsc{\chaptertitlename}}%
        }
        \hspace{10pt}%
        {\setlength\fboxsep{0pt}%
            \colorbox{myred}{\parbox[c][3cm][c]{2.5cm}{%
                \centering\color{white}\fontsize{80}{90}\selectfont\thechapter}%
            }
        }
    }
    {10pt}
    {\titlerule[2.5pt]\vskip3pt\titlerule\vskip4pt\LARGE\sffamily}

\begin{document} % Start your document

%%%%%%%%%%%% BEGIN TITLE PAGE %%%%%%%%%%%%

\thispagestyle{empty} % For the title page, no header / footer

\noindent
    \begin{minipage}{0.1\textwidth}
    \includegraphics[height=4.5em]{essex.png}
    \end{minipage}
    \hfill
    \begin{minipage}{0.5\textwidth}
    \begin{center}
        \renewcommand\familydefault{\sfdefault}
        \fontfamily{phv}\selectfont
        {\large School of Mathematics, Statistics and Actuarial Science}
    \end{center}
    \end{minipage}

\begin{center}
    \noindent\textcolor{myred}{\rule{\linewidth}{4.8pt}}
    
    \vspace{2em}
    \noindent {\LARGE \sc MA981 Dissertation}
    
    \vspace{3em}
    \noindent {\Huge{\color{myblue} YOUR PROJECT TITLE HERE}}
    
    \vspace{3em}
    \noindent {\Large \bf YOUR NAME HERE}
    \vfill
    \noindent {\Large {Supervisor:} {\color{mypurple} \bf YOUR SUPERVISOR NAME HERE}}
    
    \vspace{0.5em}
    \noindent\textcolor{myred}{\rule{\linewidth}{4.8pt}}
    
    \vspace{2em}
    {\Large \today }
    
    {\Large Colchester}
\end{center}

\clearpage

%%%%%%%%%%%% END TITLE PAGE %%%%%%%%%%%%

\tableofcontents

% If you have lots of figures with captions / numbers, uncomment the following line
% \listoffigures

% If you have lots of tables and want a list of them, uncomment the following line
% \listoftables

\chapter{Introduction}\label{ch:1}

The introduction will usually contain an overview of what is in your project document. Typically, it will be the last section you write.

\begin{theorem} \label{example-theorem}
Sometimes, you will want to state the main results of your document in the introduction.
\end{theorem}

\begin{remark}
LaTeX is clever, and automatically generates numbers for theorems, remarks and anything else you might want to label. You can give these an invisible name using \verb!\label{your-key}! and referring back to it later using \verb!\ref{your-key}!, for example the following number will be the same as the theorem above: Theorem \ref{example-theorem}.
\end{remark}

Similarly, you will want to reference external sources as you write your document. The basic way to do this is to add \verb!\bibitem{your-chosen-key}!s at the end of your document (this template has three examples), and use \verb!\cite{your-chosen-key}! to refer to it. For instance, if I wanted to cite the example document by Noether, I can write \cite{myFirstReference}.

Mathematics is added using dollar signs for in-line math, i.e. $x^2 + y^2 = z^2$, or by using open-bracket close-bracket for a displayed equation.
\[ c^2 = a^2 + b^2 - 2ab\cos \theta. \]

Ordered lists are written using the \verb!enumerate! environment:
\begin{enumerate}
    \item Hello.
    \item This is the second item in my list.
\end{enumerate}

I can also write unordered lists using \verb!itemize!:
\begin{itemize}
    \item Hello.
    \item This is now the second item in my list.
\end{itemize}

You can make figures from files as you can see in Figure \ref{fig:A3.1}. For this you need to use include graphics. 

\begin{figure}[htb]
\centerline{\includegraphics[width=1\textwidth]{Figure}}
\caption{The Gauss map $\mathbf{g}_{K}$ takes $x\in \partial K$ to the outer normal $n_x\in\mathbb{S}^{n-1}$ at that point}
\label{fig:A3.1}
\end{figure}

While writing be clear and precise and give references whenever necessary. You may like to use theorem, definition, lemma, and example environments provided by \LaTeX. For example,

Pioneering work of Emmy Noether \cite{myFirstReference} provides a connection between symmetries and conservation laws. This result, known as Noether's theorem states that

\begin{theorem}[Noether, \cite{myFirstReference}]
Every differentiable symmetry of the action of a physical system has a corresponding conservation law.
\end{theorem}

\begin{example}
This is an example.
\end{example}

\begin{lemma}
This is a lemma.
\end{lemma}

\begin{definition}
In 1950, Alan Turing published an article \cite{bookByTuring} in \textit{Mind} 
titled ``Computing Machinery and Intelligence'' where he considered the question ``Can machines think?''. This is known as {\bf Turing's Test}. 
\end{definition}

\begin{remark}
This is a very important remark. 
\end{remark}

You can also make figures  using \LaTeX packages for figures (e.g. the {\tt TikZ package}) as you can see in Figure \ref{fig:tikz}.

\begin{figure}[htb]
\begin{center}
\begin{tikzpicture}[scale=.6]
\draw[fill=purple, opacity=.3] (-3,-3) rectangle (3,3);
\fill (0,0) circle [radius=2pt];
\node at (0,4) {$A$};
\node at (4,0) {\Large $+$};
\draw[fill=red] (6,0) circle (.8cm);
\fill (6,0) circle [radius=2pt];
\node[above] at (6,1) {$\epsilon B$}; 
\node[below] at (6,-1) {\tiny radius $=\epsilon$};
\node at (8,0) {\Large $=$};

\draw [decorate,decoration={brace,amplitude=5pt, mirror}]
(-3,-3) -- (3,-3) node[below] [black,midway, yshift=-4pt] {\tiny length $=l$};

\draw[dotted, fill=blue, opacity=.1] (10,-3) rectangle (16,3);

\node at (13,0) {$A +  \epsilon B$};


\draw[dotted, fill=blue, opacity=.1] (10,-3) rectangle (16,3);

\draw[dotted, fill=red] (10,-3)--(10,-4) arc (270:180:1);
\draw[dotted, fill=red] (10,3)--(10,4) arc (90:180:1);
\draw[dotted, fill=red] (16,3)--(16,4) arc (90:0:1);
\draw[dotted, fill=red] (16,-3)--(16,-4) arc (270:360:1);

\draw[dotted] (10,-3)--(9,-3);
\draw[dotted] (16,-3)--(17,-3);

\draw[dotted] (16,3)--(16,4);
\draw[dotted] (16,3)--(17,3);

\draw[dotted] (10,3)--(9,3);


\draw[dotted, fill=purple, opacity=.1] (10,-3) rectangle (16,3);
\draw[fill=purple, opacity=.1] (16,4)--(10,4) arc (90:180:1)--(9,-3)--(9,-3) arc (180:270:1)--(10,-4)--(16,-4)--(16,-4) arc (270:360:1)--(17,-3)--(17,3) arc (0:90:1);
\end{tikzpicture}
\end{center}
\caption{Minkowski sum of a square and ball with radius $\epsilon$}
\label{fig:tikz}
\end{figure}


\chapter{Your first main chapter}\label{ch:2}

The text goes here ... 

\section{Your first section of the first main chapter}\label{sec:2.1}

... goes here.

\section{Your second section of the first main chapter}\label{sec:2.2}

... goes here.

\chapter{Your second main chapter}\label{ch:3}

The text goes here ... 

\section{Your first section of the second main chapter}\label{sec:3.1}

... goes here.

\section{Your second section of the second main chapter}\label{sec:3.2}

... goes here.

\chapter{Conclusions}\label{ch:concl}
And here is the final chapter showing how clever you are ....




% Comment the following THREE lines if you do NOT have an Appendix
\appendix
\chapter{A Long Proof}

Text goes here

% If you need more than one appendix, then just use another \chapter command
%\chapter{Yet Another Appendix}

\chapter{Another Appendix}

Text goes here




\begin{thebibliography}{999} % The '999' here tells LaTeX how much space to reserve for the bibliography item numbers. It has nothing to do with the actual number of references.

%%% Bibliography items should be below this here %%%

\bibitem{myFirstReference} % This is a key, so you can cite this with \cite{myFirstReference}
E.~Noether.
\newblock Invariante {V}ariationsprobleme.
\newblock {\em Nachr. d. K{\"o}nig. Gesellsch. d. Wiss. zu G{\"o}ttingen,
  Math-phys. Klasse, {Seite 235-157}}, 1918.

\bibitem{bookByTuring}
A.~M. Turing.
\newblock Computing machinery and intelligence.
\newblock {\em Mind}, 59:433--460, 1950.

\bibitem{myFakeBook}
J.~Fakename.
\newblock Name of book or article goes here.
\newblock \emph{Journal name}, page numbers, year, other specific info.

\end{thebibliography}

\end{document}



\end{document}